\section*{Chapitre \ref{ch:ml:targets-labels-and-sample-weights} --- Cibles, étiquettes et poids d'échantillon}

\begin{solution}{exo:ml:targets-labels-and-sample-weights:1}
Environ $1\,000/5 = 200$ étiquettes indépendantes (la \cref{prop:ml:labels:sum})~; unicité $1/5 = 0.2$ loin des extrémités.
\end{solution}

\begin{solution}{exo:ml:targets-labels-and-sample-weights:2}
Barres 1 à 8~: $c = 1, 2, 2, 2, 0, 0, 1, 1$. Unicité~: $(1 + \tfrac12)/2 = 0.75$~; $(\tfrac12 + \tfrac12 + \tfrac12)/3 =
0.5$~; $(\tfrac12 + \tfrac12)/2 = 0.5$~; $1$.
\end{solution}

\begin{solution}{exo:ml:targets-labels-and-sample-weights:3}
$1.5\times1.2~\%\times\sqrt{16} = 7.2~\%$ de chaque côté.
\end{solution}

\begin{solution}{exo:ml:targets-labels-and-sample-weights:4}
Les transactions ouvertes chaque jour avec des durées de 20~jours ont une unicité d'environ $1/20$, si bien que le nombre effectif est d'environ 1\,250 et la
$t$-statistique d'environ $6.1/\sqrt{20} = 1.36$~: non significative.
\end{solution}

\begin{solution}{exo:ml:targets-labels-and-sample-weights:5}
La $t$-statistique est la moyenne divisée par l'écart type, multipliée par la racine carrée du nombre effectif de transactions.
En conservant un tiers des transactions, on augmente la moyenne de moitié (de 0.064 à 0.097 largeurs), mais on divise le nombre effectif par environ
trois, si bien que $t$ baisse (de 3.48 à 3.19). Elle augmenterait si le modèle secondaire trouvait des transactions d'espérance négative à écarter
(une information que le modèle primaire ne possède pas), et non simplement des transactions d'espérance positive plus faible.
\end{solution}

\begin{solution}{exo:ml:targets-labels-and-sample-weights:6}
L'écart type calculé sur l'ensemble de l'échantillon utilise le futur (biais d'anticipation) et ignore le regroupement de volatilité~: les barrières sont trop
larges dans les périodes calmes et trop étroites en période de stress, si bien que les étiquettes se règlent pour des raisons étrangères à la transaction. Il faut utiliser la volatilité
connue à la date de l'événement (EWMA ou GARCH, livre~4, chapitre~18).
\end{solution}

\begin{solution}{exo:ml:targets-labels-and-sample-weights:7}
\texttt{label\_stats(k=2.0)}~: profit 4.0~\%, stop 3.3~\%, verticale 92.7~\%~; les étiquettes durent 9.8 jours, unicité 0.1025
(0.123 pour $k = 1$). Des barrières plus larges sont rarement touchées dans les dix jours, si bien que l'\omterm{def:ml:targets-labels-and-sample-weights:triple}{étiquette à triple barrière} converge vers
l'\omterm{def:ml:targets-labels-and-sample-weights:fixed}{étiquette à horizon fixe}.
\end{solution}

\begin{solution}{exo:ml:targets-labels-and-sample-weights:8}
Avec une unicité constante, l'unicité cumulée est proportionnelle au rang de l'étiquette~: l'étiquette à mi-parcours a le poids
$0.5 + 0.5\times0.5 = 0.75$, et la moyenne d'une rampe linéaire de 0.5 à 1 vaut 0.75. Avec une unicité variable, les périodes encombrées
(faible unicité) font avancer la rampe lentement~; leurs nombreuses étiquettes partagent donc des poids voisins, et la décroissance suit
l'information, et non le calendrier~; la moyenne dépend alors de l'endroit où tombent les périodes encombrées.
\end{solution}

\begin{solution}{pb:ml:targets-labels-and-sample-weights:1}
\begin{enumerate}
\item Profit 22.7~\%, stop 20.3~\%, verticale 57.0~\%.
\item 8.3 jours~: une étiquette se règle au premier contact, avant le dixième jour dans 43~\% des événements.
\item Horizon fixe~: haussière (un gagnant, $+1.84~\%$)~; triple barrière~: stoppée à $-5.04~\%$ le troisième jour, un perdant. L'\omterm{def:ml:targets-labels-and-sample-weights:triple}{étiquette à triple barrière}
décrit la transaction.
\item 1.9~\%~; les tailles diffèrent pour chaque étiquette touchée, puisque les rendements à barrières sont tronqués comme ceux de la transaction.
\item Concurrence 10, unicité 0.100.
\item $20\times2\,458/10 = 4\,916$~; \omterm{def:ml:targets-labels-and-sample-weights:triple}{étiquettes à triple barrière} 6\,040.
\item 9.4 sur les lignes, 3.3 ajustée~; rapport $2.85 = \sqrt{1/0.123}$ (la \cref{prop:ml:labels:sum}, partie 2).
\item Erreurs types, $t$-statistiques, intervalles de bootstrap, scores hors sac et tout test de significativité qui compte
les lignes.
\item 0.65 et 0.69.
\item Simple 0.504, 0.517, 0.509~; petit 0.511, 0.512, 0.506~; pondéré 0.506, 0.511, 0.510.
\item Les arbres étaient régularisés et aucune variable ne permet d'identifier la date, si bien que des quasi-doublons dans un échantillon bootstrap
ne pouvaient pas être mémorisés en un avantage.
\item Arbres profonds avec une variable à évolution lente qui repère le temps (une variable de régime, une tendance du niveau)~: des voisins qui se chevauchent
partagent alors à la fois les variables et les étiquettes, et des échantillons bootstrap de pleine taille les ajustent.
\item Taux de réussite de 54.5~\% contre 53.0~\% et 0.097 contre 0.064 largeurs par transaction, sur un tiers des transactions~; la
$t$-statistique ajustée du chevauchement passe de 3.48 à 3.19.
\item Il choisit lesquelles des transactions du modèle primaire prendre et pour quelle taille~; le sens, là où se trouve le signal, vient du
modèle primaire.
\item Oui, proportionnellement à $2p - 1$ ou à une fraction de type Kelly (chapitre~11), qui conserve chaque transaction mais la dimensionne~:
le filtrage est le cas particulier des tailles 0 et 1.
\item \emph{Résultat nommé}~: 4\,916 étiquettes effectives sur 48\,980 (10~\%)~; $t$ de 9.4 sur les lignes, 3.3 ajustée~; la pondération
par unicité~: gain d'AUC par rapport au bagging simple de $-0.001$ en moyenne, dans $\pm0.007$ selon le marché.
\item Environ $20\times2\,450/20 = 2\,450$~: deux fois moins.
\item Lignes, durées ($t_0$, $t_1$), somme des unicités et nombre de dates indépendantes.
\item Les deux proviennent des durées~: la purge retire les étiquettes d'entraînement dont les durées chevauchent le pli de test~; l'unicité mesure
le degré de chevauchement des durées au sein d'un échantillon.
\item Une étiquette est le règlement d'une transaction~: ce que la position a rapporté, de l'instant de décision à la sortie.
\end{enumerate}
\end{solution}

\begin{solution}{iq:ml:targets-labels-and-sample-weights:1}
Elle ignore le chemin (stops et objectifs), ignore la volatilité (un seuil de rendement fixe n'a pas le même sens selon les régimes), et
des étiquettes consécutives se chevauchent sur 19 jours sur 20, si bien que l'échantillon vaut environ un vingtième de ses
lignes.
\iqlookfor{chemin, volatilité et chevauchement.}
\end{solution}

\begin{solution}{iq:ml:targets-labels-and-sample-weights:2}
Pour chaque événement, on fixe une barrière de profit et un stop en multiples de la volatilité courante, ainsi qu'une barrière verticale à une
durée de détention maximale~; l'étiquette est le résultat de la première barrière touchée, et elle enregistre la date de son règlement.
\iqlookfor{mise à l'échelle par la volatilité, premier contact et date de règlement conservée pour la purge.}
\end{solution}

\begin{solution}{iq:ml:targets-labels-and-sample-weights:3}
Par action, environ $2\,520/10 = 252$ étiquettes indépendantes, soit environ 100\,000 au total, et moins puisque les actions sont corrélées
(un facteur commun fait qu'une date vaut quelques actions indépendantes). L'échantillon effectif favorise les modèles régularisés à faible variance
et des erreurs types honnêtes.
\iqlookfor{unicité multipliée par les lignes, puis corrélation en coupe.}
\end{solution}

\begin{solution}{iq:ml:targets-labels-and-sample-weights:4}
Un classifieur secondaire estime la probabilité que l'appel d'un modèle primaire soit juste, entraîné sur des étiquettes indiquant si chaque appel
a rapporté de l'argent~; sa sortie filtre ou dimensionne les transactions. Utile lorsque le sens primaire provient d'une source que le
modèle ne peut pas reproduire (un avis discrétionnaire, une règle) et que le modèle secondaire voit quand elle fonctionne.
\iqlookfor{sens contre taille, et le fait qu'il ne peut pas ajouter un signal que le modèle primaire ne possède pas.}
\end{solution}

\begin{solution}{iq:ml:targets-labels-and-sample-weights:5}
Les observations hors sac sont celles qu'un arbre n'a pas tirées, mais, les étiquettes se chevauchant, leurs voisines (partageant l'essentiel
de l'étiquette et des variables similaires) ont été tirées~: le score n'est pas hors échantillon. La période de test ultérieure ne présente pas
de tels voisins. Il faut plutôt utiliser une validation purgée et ordonnée dans le temps.
\iqlookfor{la fuite du chevauchement vers l'estimation hors sac.}
\end{solution}

\begin{solution}{iq:ml:targets-labels-and-sample-weights:6}
Chaque innovation entre dans $h$ des $n$ sommes, si bien que la moyenne des sommes vaut environ $\frac hn\sum_t\varepsilon_t$, de variance
$h^2\sigma^2/n$~; en traitant les sommes comme indépendantes, on obtient $\Var(h\text{-somme})/n = h\sigma^2/n$. Le rapport vaut $h$.
\iqlookfor{l'argument de comptage et sa conséquence pour les $t$-statistiques ($\sqrt h$).}
\end{solution}
